\label{bkm:Ref160277935}\index{BPatch_flowGraph: : : : : : :!}{

The BPatch_flowGraph class represents the control flow graph of a function. It provides methods for discovering the

basic blocks and loops within the function (using which a caller can navigate the graph). A BPatch_flowGraph object

can be obtained by calling the getCFG method of a BPatch_function object.}



{

bool containsDynamicCallsites()}



{

Return true if the control flow graph contains any dynamic call sites (e.g., calls through a function pointer).}



{

void getAllBasicBlocks\index{getAllBasicBlocks: : : : : : :!}(std::set<BPatch_basicBlock*>&)



{

void getAllBasicBlocks\index{getAllBasicBlocks: : : : : : :!}(BPatch_Set<BPatch_basicBlock*>&)

}



{

Fill the given set with pointers to all basic blocks in the control flow graph. BPatch_basicBlock is described in

section \ref{bkm:Ref161214701}.}



{

void getEntryBasicBlock\index{getEntryBasicBlock: : : : : :

:!}(std::vector<BPatch_basicBlock*>&)}



{

Fill the given vector with pointers to all basic blocks that are entry points to the function. BPatch_basicBlock is

described in section \ref{bkm:Ref161214701}.}



{

void getExitBasicBlock\index{getExitBasicBlock: : : : : :

:!}(std::vector<BPatch_basicBlock*>&)}



{

Fill the given vector with pointers to all basic blocks that are exit points of the function. BPatch_basicBlock is

described in section \ref{bkm:Ref161214701}.}



{

void getLoops\index{getLoops: : : : : : :!}(std::vector<BPatch_basicBlockLoop*>&)}



{

Fill the given vector with a list of all natural (single entry) loops in the control flow graph.}





\bigskip



{

void getOuterLoops\index{getLoops: : : : : : :!}(std::vector<BPatch_basicBlockLoop*>&)}



{

Fill the given vector with a list of all natural (single entry) outer loops in the control flow graph.}



{

BPatch_loopTreeNode *getLoopTree\index{getLoops: : : : : : :!}()}



{

Return the root node of the tree of loops in this flow graph.}





\bigskip



{

enum BPatch_procedureLocation { BPatch_locLoopEntry, BPatch_locLoopExit, BPatch_locLoopStartIter,

BPatch_locLoopEndIter }}





\bigskip



{

std::vector<BPatch_point*> *findLoopInstPoints\index{getLoops: : : : : : :!}(const

BPatch_procedureLocation loc, BPatch_basicBlockLoop *loop);}



{

Find instrumentation points for the given loop that correspond to the given location: loop entry, loop exit, the start

of a loop iteration and the end of a loop iteration. BPatch_locLoopEntry and BPatch_locLoopExit instrumentation

points respectively execute once before the first iteration of a loop and after the last iteration.

BPatch_locLoopStartIter and BPatch_locLoopEndIter respectively execute at the beginning and end of each loop

iteration.}



{

BPatch_basicBlock* findBlockByAddr(Dyninst::Address addr);}



{

Find the basic block within this flow graph that contains addr. Returns NULL on failure. This method is inefficient but

guaranteed to succeed if addr is present in any block in this CFG.}





\bigskip



{

[NOTE: Dyninst is not always able to generate a correct flow graph in the presence of indirect jumps. If a function

has a case statement or indirect jump instructions, the targets of the jumps are found by searching instruction

patterns (peep-hole). The instruction patterns generated are compiler specific and the control flow graph analyses

include only the ones we have seen. During the control flow graph generation, if a pattern that is not handled is used

for case statement or multi-jump instructions in the function address space, the generated control flow graph may not

be complete.]}



